\paragraph{An original-action infinite-volume ground representation at fixed spacing.}
The full three-dimensional Weyl/contact action now admits a
translation-invariant, gauge-invariant ground state at its original
zero chemical potential, without choosing a particle-number sector.
This is a compactness existence theorem, not a computed or natively
prepared state.  Keep fixed finite $c>0$, $U>0$ and the complete
quasi-local CAR algebra with two modes per site.  The previously proved
full-Fock locality estimate supplies the strongly continuous dynamics
$\tau_t$.  On each even-sided cube $\Lambda_L$ of side $2L$, take the
normalized projector onto the absolute minimum of $H_{\Lambda_L}$,
extend this even state by the exterior vacuum, and average its translates
over all sites of $\Lambda_L$.  A diagonal subsequence of the bounded
local matrix entries defines a state $\omega$; positivity, normalization
and CAR compatibility pass to the limit.  The vanishing
boundary-to-volume ratio makes $\omega$ translation invariant.

For a local $A$, including an odd charged probe, the local derivation
$[H,A]$ contains only interaction terms touching its support.  Whenever
the support and its nearest neighbors lie inside the box,
\[
 \omega_L\!\left(A^*[H,A]\right)
 =\operatorname{Tr}\rho_L A^*(H_{\Lambda_L}-E_L)A\ge0.
\]
Bad translates occupy a vanishing fraction.  Their real-part lower
defect and imaginary magnitude are each bounded by
$2\|A\|^2|\operatorname{supp}A|(U+6c)$, so
$\omega(A^*[H,A])\ge0$ in the limit.  This uses the full ground
minimum and would fail for a merely canonical ground with an
unrestricted charged probe.

For a radius-$R$ buffer, let $B_R(t)$ denote the existing
volume-uniform local-dynamics boundary bound.  Comparing both finite
and infinite dynamics to this intermediate cube gives, for $\|A\|\le1$,
\[
 \bigl|\nu_L(\tau_t(A))-\nu_L(A)\bigr|
 \le \min\!\left\{2,\,2B_R(t)
       +2\left[1-(1-R/L)^3\right]\right\}.
\]
First $L\to\infty$, then $R\to\infty$, establishes stationarity.
The same comparison applies to two-point correlations.  Every finite
full-ground correlation has a positive spectral measure supported on
$[0,\infty)$; compact-time convergence and bounded correlations pass
all Schwartz Fourier tests supported on negative energies to the limit.
The GNS implementation therefore satisfies
\[
 U_t\pi(A)\Omega=\pi(\tau_t(A))\Omega,
 \qquad U_t=e^{itL_\omega},\qquad
 L_\omega\ge0,\quad L_\omega\Omega=0.
\]
The generator measures excitation energy in this representation;
no chemical potential has been inserted into the source action.

The state can be chosen at nonzero density.  Tiling each even box by
the already certified four-particle cube vector gives contact zero,
energy $-2\sqrt3 c$ per eight-site tile, and zero expectation for every
cross-tile hopping.  Hence $E_L/|\Lambda_L|\le-\sqrt3 c/4$.
Cartesian difference operators commute and the Pauli mixed terms
cancel, giving $h^2\le3c^2I$ and $H\ge-\sqrt3 cN$.
Thus every such full ground has density at least $1/4$.
The exact real-bipartite particle--hole identity
$E(2V-N)=E(N)+U(V-N)$ excludes $N>V$ from an absolute minimum when
$U>0$.  Spatial averaging before the subsequence passes both bounds:
\[
 \frac14\le\omega(n_{0\uparrow}+n_{0\downarrow})\le1.
\]
Half-filling and a particular density are not assumed or concluded.
The original ground energy per site also has a limit $e_\infty$.
For every even cube side $\ell$, tiling larger boxes gives the
computable bracket, conditional on its finite full-ground energy,
\[
 \frac{E_\ell}{\ell^3}-\frac{3c}{\ell}
 \le e_\infty\le\frac{E_\ell}{\ell^3},
 \qquad -\sqrt3 c\le e_\infty\le-\frac{\sqrt3 c}{4}.
\]
The lower surface term counts intertile hopping norms; a product of
gauge-invariant tile grounds gives the upper bound.  Untiled boundary
terms vanish per volume, and the bracket width tends to zero.
For the constructed averaged ground, $\omega(h_0)=e_\infty$, where
$h_0$ assigns half of each incident bond to the origin.  This supplies
an energy-density observable, not a state-convergence rate or rest mass.

The contact interaction remains represented:
\[
 [H_U-H_0,a_0^*a_2]
   =U a_0^*a_2(n_1-n_3)\ne0.
\]
Every unital GNS representation is injective on each full local matrix
algebra, so this quartic operator cannot disappear in representation.
This faithfulness of the representation does not assert faithfulness
of the ground-state functional.  The claim is a nonzero represented
interaction derivation, not a scattering or connected-vacuum-correlator
theorem.

An exact source countercontrol prevents simply substituting the free
filled sea.  For a pure Slater projector with $P_{xx}=I_2/2$,
\[
 \operatorname{Var}(D)
   =\frac V{16}+\sum_{x\ne y}|\det P_{xy}|^2\ge\frac V{16}.
\]
If it is also a kinetic eigenstate, its interacting energy variance is
$U^2\operatorname{Var}(D)>0$.  On the original cube, the complete
3844-term negative-band filling has $N=8$, $\langle D\rangle=2$ and
$\operatorname{Var}(H)=2/3$ at $U=1$.  The exact checks and proof are
retained in \texttt{model/thermodynamic\_ground/results.json}.
For local CAR probes, the same subsequence also supplies the analytic
retarded response
\[
 G_{AB}(z)=-i\int_0^\infty e^{izt}
       \omega(\{\tau_t(A),B^*\})\,dt,\qquad \operatorname{Im}z>0.
\]
Each time tail after $T$ is at most
$2\|A\|\|B\|e^{-\eta T}/\eta$ for $\operatorname{Im}z\ge\eta>0$.
Compact-time locality and weak-star convergence, followed by this tail
bound, give uniform convergence on compact subsets of the upper
half-plane.  The unknown finite-state convergence error is retained;
no real-axis pole, local inverse or metric is inferred.
The result closes fixed-spacing ground-state existence and the
positive-energy representation for the actual interacting lattice.
It supplies no selected ground functional, convergence rate, native
preparation, repair stationarity, uniform gap, Lorentz continuum,
physical particle assignment or Einstein dynamics.
